French uses a vigesimal (base-20) counting system for numbers 70--99, where expressions like \emph{quatre-vingt-dix-sept} (literally ``four-twenty-ten-seven'') encode implicit arithmetic operations ($4 \times 20 + 10 + 7 = 97$).
We investigate how large language models represent and process these linguistically complex numerals through two complementary analyses: behavioral evaluation of \gptmodel on number conversion, arithmetic, and counting tasks, and representational analysis of \mistral's internal hidden states across standard French, Belgian French, and digit forms.
At the behavioral level, \gptmodel achieves 100\% accuracy on all tasks, including vigesimal numbers---refuting the hypothesis that irregular number words cause systematic errors in frontier models.
However, probing \mistral's representations reveals that the vigesimal system leaves a distinctive internal fingerprint: 70\% of vigesimal numbers' nearest neighbors in representation space share a linguistic base (e.g., the \quatrevingt prefix) rather than numeric proximity.
French number words are more linearly decodable ($R^2 = 0.979$) than digit strings ($R^2 = 0.943$), and Belgian French's regular decimal system produces representations with the highest ordinal correlation ($\rho = 0.591$ vs.\ $0.520$ for standard French).
These findings demonstrate that models ``think in vigesimal'' even when they ``answer in decimal,'' revealing how linguistic structure shapes numerical representations beyond what behavioral benchmarks can detect.
